% GENERATED FILE. Edit canonical lessons, then run npm run generate:books.
% canonical-source-hash: fnv1a64:8db4e7454219d058
% canonical-lessons: BN-C476-dom, BN-C476-hechki, BN-C476-dhekur, BN-C476-hapano, BN-C476-katrano

\chapter{Out of Breath}
\label{ch:bn-out-of-breath}
\hlchaptermodality{476}

\begin{chapteropening}
\textbf{By the end of this chapter:} \emph{I can say that I am out of breath, have hiccups or am groaning.}

Complete the last lesson of chapter 476: I can say that I am out of breath, have hiccups or am groaning.
\end{chapteropening}

\section[\texorpdfstring{\bn{দম} (dôm)}{দম (dôm)}]{\bn{দম} (dôm) — breath; stamina}

\label{lesson:BN-C476-dom}



\begin{quote}
Before the new one: say the Bengali for to swell, then the Bengali for to stumble.
\end{quote}

\subsection*{You'll want to know: \bn{দম}}
\textbf{\bn{দম}} — \emph{dôm} — \textquotedblleft{}breath; stamina\textquotedblright{}.

\begin{quote}
\bn{সাঁতার} \bn{কাটলে} \bn{দম} \bn{বাড়ে।}
\end{quote}

\begin{quote}
\emph{sā̃tār kāṭle dôm bāṛe.}
\end{quote}

\begin{quote}
Swimming builds up your stamina.
\end{quote}

\bn{দম} \bn{নেওয়া}, to take a breather.

\begin{cousinweb}
\textbf{\bn{দম}} (\emph{dôm}) is usually traced to Persian \emph{dam}, \textquotedblleft{}breath\textquotedblright{}.
\end{cousinweb}

\begin{grammarlens}[title={What you've built}]
One more word for describing what you lived and felt.
\end{grammarlens}

\subsection*{Your turn}
\begin{itemize}
\raggedright
  \item \emph{Say it:} \emph{dôm}
  \item \emph{Say it:} \emph{dôm}, once more
  \item \emph{Say it:} \emph{sā̃tār kāṭle dôm bāṛe.} — Swimming builds up your stamina.
  \item \emph{Recall:} say the Bengali for to swell, then the Bengali for to stumble, then say \emph{dôm} again
\end{itemize}

\begin{tcolorbox}[breakable,colback=teal!4,colframe=teal!35!black,title={Before you move on}]
What does \bn{দম} mean? (\textquotedblleft{}Breath; stamina\textquotedblright{}.) Say it once more.
\end{tcolorbox}
\section[\texorpdfstring{\bn{হেঁচকি} (hẽchki)}{হেঁচকি (hẽchki)}]{\bn{হেঁচকি} (hẽchki) — a hiccup; hiccups}

\label{lesson:BN-C476-hechki}



\begin{quote}
Before the new one: say the Bengali for to stumble, then the Bengali for stamina.
\end{quote}

\subsection*{You'll want to know: \bn{হেঁচকি}}
\textbf{\bn{হেঁচকি}} — \emph{hẽchki} — \textquotedblleft{}a hiccup; hiccups\textquotedblright{}.

\begin{quote}
\bn{তাড়াতাড়ি} \bn{খেলে} \bn{আমার} \bn{হেঁচকি} \bn{হয়।}
\end{quote}

\begin{quote}
\emph{tāṛātāṛi khele āmār hẽchki hôy.}
\end{quote}

\begin{quote}
If I eat too fast I get hiccups.
\end{quote}

\begin{cousinweb}
\textbf{\bn{হেঁচকি}} (\emph{hẽchki}) is an imitative word, its sound standing for the catch of a hiccup. No dictionary we use records an older form, so its history is not settled.
\end{cousinweb}

\begin{grammarlens}[title={What you've built}]
One more word for describing what you lived and felt.
\end{grammarlens}

\subsection*{Your turn}
\begin{itemize}
\raggedright
  \item \emph{Say it:} \emph{hẽchki}
  \item \emph{Say it:} \emph{hẽchki}, once more
  \item \emph{Say it:} \emph{tāṛātāṛi khele āmār hẽchki hôy.} — If I eat too fast I get hiccups.
  \item \emph{Recall:} say the Bengali for to stumble, then the Bengali for stamina, then say \emph{hẽchki} again
\end{itemize}

\begin{tcolorbox}[breakable,colback=teal!4,colframe=teal!35!black,title={Before you move on}]
What does \bn{হেঁচকি} mean? (\textquotedblleft{}A hiccup; hiccups\textquotedblright{}.) Say it once more.
\end{tcolorbox}
\section[\texorpdfstring{\bn{ঢেকুর} (ḍhekur)}{ঢেকুর (ḍhekur)}]{\bn{ঢেকুর} (ḍhekur) — a burp; a belch}

\label{lesson:BN-C476-dhekur}



\begin{quote}
Before the new one: say the Bengali for stamina, then the Bengali for hiccups.
\end{quote}

\subsection*{You'll want to know: \bn{ঢেকুর}}
\textbf{\bn{ঢেকুর}} — \emph{ḍhekur} — \textquotedblleft{}a burp; a belch\textquotedblright{}.

\begin{quote}
\bn{ভারী} \bn{খাবারের} \bn{পরে} \bn{দাদু} \bn{ঢেকুর} \bn{তুললেন।}
\end{quote}

\begin{quote}
\emph{bhārī khābārer pôre dādu ḍhekur tullen.}
\end{quote}

\begin{quote}
After the heavy meal my grandfather burped.
\end{quote}

\bn{ঢেকুর} \bn{তোলা}, to burp.

\begin{cousinweb}
\textbf{\bn{ঢেকুর}} (\emph{ḍhekur}) is an imitative word for a belch; it shows the new letter \bn{ঢ}. No dictionary we use records an older form, so its history is not settled.
\end{cousinweb}

\begin{grammarlens}[title={What you've built}]
One more word for describing what you lived and felt.
\end{grammarlens}

\subsection*{Your turn}
\begin{itemize}
\raggedright
  \item \emph{Say it:} \emph{ḍhekur}
  \item \emph{Say it:} \emph{ḍhekur}, once more
  \item \emph{Say it:} \emph{bhārī khābārer pôre dādu ḍhekur tullen.} — After the heavy meal my grandfather burped.
  \item \emph{Recall:} say the Bengali for stamina, then the Bengali for hiccups, then say \emph{ḍhekur} again
\end{itemize}

\begin{tcolorbox}[breakable,colback=teal!4,colframe=teal!35!black,title={Before you move on}]
What does \bn{ঢেকুর} mean? (\textquotedblleft{}A burp; a belch\textquotedblright{}.) Say it once more.
\end{tcolorbox}
\section[\texorpdfstring{\bn{হাঁপানো} (hā̃pāno)}{হাঁপানো (hā̃pāno)}]{\bn{হাঁপানো} (hā̃pāno) — to pant; to be out of breath}

\label{lesson:BN-C476-hapano}



\begin{quote}
Before the new one: say the Bengali for hiccups, then the Bengali for a burp.
\end{quote}

\subsection*{You'll want to know: \bn{হাঁপানো}}
\textbf{\bn{হাঁপানো}} — \emph{hā̃pāno} — \textquotedblleft{}to pant; to be out of breath\textquotedblright{}.

\begin{quote}
\bn{তাড়াতাড়ি} \bn{হেঁটে} \bn{এসে} \bn{দাদু} \bn{হাঁপাচ্ছিলেন।}
\end{quote}

\begin{quote}
\emph{tāṛātāṛi hẽṭe ese dādu hā̃pāchchhilen.}
\end{quote}

\begin{quote}
Grandfather was panting after walking fast.
\end{quote}

Past (he/she …): \bn{তিনি} \bn{হাঁপালেন}, \emph{tini hā̃pālen}.

\begin{cousinweb}
\textbf{\bn{হাঁপানো}} (\emph{hā̃pāno}) is listed by CDIAL for Bengali under a reconstructed root for panting (CDIAL 13973); \bn{হাঁপানি}, asthma, is its noun.
\end{cousinweb}

\begin{grammarlens}[title={What you've built}]
One more word for describing what you lived and felt.
\end{grammarlens}

\subsection*{Your turn}
\begin{itemize}
\raggedright
  \item \emph{Say it:} \emph{hā̃pāno}
  \item \emph{Say it:} \emph{hā̃pāno}, once more
  \item \emph{Say it:} \emph{tāṛātāṛi hẽṭe ese dādu hā̃pāchchhilen.} — Grandfather was panting after walking fast.
  \item \emph{Recall:} say the Bengali for hiccups, then the Bengali for a burp, then say \emph{hā̃pāno} again
\end{itemize}

\begin{tcolorbox}[breakable,colback=teal!4,colframe=teal!35!black,title={Before you move on}]
What does \bn{হাঁপানো} mean? (\textquotedblleft{}To pant; to be out of breath\textquotedblright{}.) Say it once more.
\end{tcolorbox}
\section[\texorpdfstring{\bn{কাতরানো} (kātrāno)}{কাতরানো (kātrāno)}]{\bn{কাতরানো} (kātrāno) — to groan (in pain); to writhe}

\label{lesson:BN-C476-katrano}



\begin{quote}
Before the new one: say the Bengali for a burp, then the Bengali for to pant.
\end{quote}

\subsection*{You'll want to know: \bn{কাতরানো}}
\textbf{\bn{কাতরানো}} — \emph{kātrāno} — \textquotedblleft{}to groan (in pain); to writhe\textquotedblright{}.

\begin{quote}
\bn{পেটব্যথায়} \bn{রোগীটা} \bn{সারারাত} \bn{কাতরাল।}
\end{quote}

\begin{quote}
\emph{peṭbæthāy rogīṭā sārārāt kātrālô.}
\end{quote}

\begin{quote}
The patient groaned all night with a stomach ache.
\end{quote}

Past (he/she …): \bn{তিনি} \bn{ব্যথায়} \bn{কাতরালেন}, \emph{tini bæthāy kātrālen}.

\begin{cousinweb}
\textbf{\bn{কাতরানো}} (\emph{kātrāno}) is the verb of \emph{kātôr}, \textquotedblleft{}distressed\textquotedblright{}, Sanskrit \emph{kātara}, \textquotedblleft{}timid, despairing, agitated\textquotedblright{} (Monier-Williams).
\end{cousinweb}

\begin{grammarlens}[title={What you've built}]
One more word for describing what you lived and felt.
\end{grammarlens}

\subsection*{Your turn}
\begin{itemize}
\raggedright
  \item \emph{Say it:} \emph{kātrāno}
  \item \emph{Say it:} \emph{kātrāno}, once more
  \item \emph{Say it:} \emph{peṭbæthāy rogīṭā sārārāt kātrālô.} — The patient groaned all night with a stomach ache.
  \item \emph{Recall:} say the Bengali for a burp, then the Bengali for to pant, then say \emph{kātrāno} again
\end{itemize}

\begin{tcolorbox}[breakable,colback=teal!4,colframe=teal!35!black,title={Before you move on}]
What does \bn{কাতরানো} mean? (\textquotedblleft{}To groan (in pain); to writhe\textquotedblright{}.) Say it once more.
\end{tcolorbox}
